\documentclass[10pt,english]{beamer}

\usetheme{default}
\setbeamertemplate{navigation symbols}{\textcolor{blue}{\insertframenumber ~/ \inserttotalframenumber}}

\usepackage{amsmath,amssymb,amsthm}
\usepackage{stmaryrd}
\usepackage{enumerate}
\usepackage{stfloats}
\usepackage{bbm}
\usepackage{pdfpages}
\usepackage{framed}

\usepackage{tikz,pgf,pgfplots}
\usepackage{algorithm,algorithmic}
\usepgflibrary{shapes}
\usetikzlibrary{%
  arrows,%
  arrows.meta,
  shapes.misc,% wg. rounded rectangle
  shapes.arrows,%
  shapes,%
  calc,%
  chains,%
  matrix,%
  positioning,% wg. " of "
  scopes,%
  decorations.pathmorphing,% /pgf/decoration/random steps | erste Graphik
  shadows,%
  backgrounds,%
  fit,%
  petri,%
  quotes
}

%\pgfplotsset{compat=1.12}

%\usetheme{Frankfurt}
%\usecolortheme{ldpc}
\useinnertheme{rounded}
\usecolortheme{whale}
\usecolortheme{orchid}


\newcommand{\ul}[1]{\underline{#1}}
\renewcommand{\Pr}{\mathbb{P}}

\newcommand{\getpdfpages}[2]{\begingroup
  \setbeamercolor{background canvas}{bg=}
  \addtocounter{framenumber}{1}
  \includepdf[pages={#1},%
  pagecommand={%
    \expandafter\def\expandafter\insertshorttitle\expandafter{%
      \insertshorttitle\hfill\insertframenumber\,/\,\inserttotalframenumber}}%
  ]{#2}
  \endgroup}

\newcommand{\backupbegin}{
   \newcounter{finalframe}
   \setcounter{finalframe}{\value{framenumber}}
}
\newcommand{\backupend}{
   \setcounter{framenumber}{\value{finalframe}}
}

 \setbeamercolor{bibliography entry author}{fg=black}
 \setbeamercolor{bibliography entry title}{fg=black}
 \setbeamercolor{bibliography entry location}{fg=black}
 \setbeamercolor{bibliography entry note}{fg=black}
 
 \setbeamerfont{bibliography item}{size=\footnotesize}
 \setbeamerfont{bibliography entry author}{size=\footnotesize}
 \setbeamerfont{bibliography entry title}{size=\footnotesize}
 \setbeamerfont{bibliography entry location}{size=\footnotesize}
 \setbeamerfont{bibliography entry note}{size=\footnotesize}
 \setbeamertemplate{bibliography item}{\insertbiblabel}
 
\newlength\tikzwidth
\newlength\tikzheight

\def\checkmark{\tikz\fill[scale=0.4](0,.35) -- (.25,0) -- (1,.7) -- (.25,.15) -- cycle;}
\def\greencheck{{\color{green}\checkmark}}
\def\scalecheck{\resizebox{\widthof{\checkmark}*\ratio{\widthof{x}}{\widthof{\normalsize x}}}{!}{\checkmark}}
\def\xmark{\tikz [x=1.4ex,y=1.4ex,line width=.2ex, red] \draw (0,0) -- (1,1) (0,1) -- (1,0);}
\def\redx{{\color{red}\xmark}}

\renewcommand{\footnotesep}{-2pt}

\newif\ifslow
\slowtrue

\newcommand{\mc}[1]{\mathcal{#1}}
\newcommand{\mbb}[1]{\mathbb{#1}}
%\newcommand{\expt}{\mbb{E}}
%\newcommand{\dd}{\mathrm{d}}
\newcommand{\Interior}[1]{\ensuremath{{#1}^{\circ}}}
\newcommand{\Closure}[1]{\ensuremath{\overline{#1}}}
\newcommand{\Complement}[1]{\ensuremath{{#1}^{c}}}

\newcommand{\Expect}{\ensuremath{\mathrm{E}}}
\newcommand{\vecnot}{\underline}
\newcommand{\RealNumbers}{\ensuremath{\mathbb{R}}}
\newcommand{\RationalNumbers}{\mathbb{Q}}
\newcommand{\ComplexNumbers}{\mathbb{C}}
\newcommand{\Real}{\mathrm{Re}}
\newcommand{\Span}{\mathrm{span}}
\newcommand{\Rank}{\mathrm{rank}}
\newcommand{\Nullity}{\mathrm{nullity}}
\newcommand{\Trace}{\mathrm{tr}}
\newcommand{\Diag}{\mathrm{diag}}
\newcommand{\dd}{\mathrm{d}}
\DeclareMathOperator*{\esssup}{ess\,sup}

% Use < , > inner product
\newcommand{\inner}[2]{{\left\langle #1 \mskip2mu , #2 \right\rangle}}
\newcommand{\tinner}[2]{{\langle #1 \mskip1mu , #2 \rangle}}

% Use < | > inner product
%\newcommand{\inner}[2]{{\left\langle #1 \mskip2mu \middle| \mskip2mu #2 \right\rangle}}
%\newcommand{\tinner}[2]{{\langle #1 \mskip1mu | \mskip1mu  #2 \rangle}}




\begin{document}

\ifslow

\title{ECE 586: Vector Space Methods \\ Chapter 2: Metric Spaces and Topology}
\author{Henry D. Pfister \\ Duke University}
\date{September 4th -- 16th, 2019}
\maketitle

\begin{frame}{Introduction}

\begin{itemize}
\setlength\itemsep{3mm}
\item<1-> What is topology and why do we study it? \vspace{1mm}
\begin{itemize} 
  \setlength\itemsep{1.5mm}
  \item<1-> \textcolor{blue}{Study of geometric properties preserved by continuous deformations}
  \item<2-> Why? Engineers approximate real things by mathematical objects
  \item<2-> Q1: Can a matrix be approximated well by a lower rank matrix?
  \item<2-> Q2: Can a function be approximated well by a degree-2 polynomial?
  \item<2-> In engineering, a topology is typically defined using a metric
\end{itemize}

\vspace{1mm}

\item<3-> Metric Spaces \vspace{1mm}
\begin{itemize} 
  \setlength\itemsep{1.5mm}
  \item A \textcolor{blue}{metric space} $(X,d)$ is a set $X$ along with a well-defined metric $d$
  \item A \textcolor{blue}{metric} on a set $X$ is a function \textcolor{blue}{$d \colon X \times X \rightarrow \mathbb{R} $} that satisfies: \vspace{1mm}
  \begin{itemize}
  \setlength\itemsep{1.5mm}
  \item $d(x,y) \geq 0 \quad \forall x, y \in X$; with equality if and only if $x = y$
  \item $d(x,y) = d(y,x) \quad \forall x, y \in X$
  \item $d(x,y) + d(y,z) \geq d(x,z) \quad \forall x, y, z \in X$.
  \end{itemize}
  \item $d(x,y)$ is called the \textcolor{blue}{distance} between \textcolor{blue}{points} $x$ and $y$
  \item Whiteboard Examples
\end{itemize}
\end{itemize}
\end{frame}

\begin{frame}{Useful Abstractions}

%\vspace*{-2mm}
\hfill
\scalebox{0.9}{%}
\begin{tikzpicture}
\draw[very thick] (0,0) rectangle (3,2.75) node[below left=0.2cm and 0.2cm] {$X$};
\draw[thick,dashed,blue,fill=blue!10] (1.25,1.25) circle (1) node[above=0.3cm] {$B(x,\epsilon)$};
\node [label=right:{$x$},draw,fill=black,circle,inner sep=0pt,minimum size=2pt] at (1.25,1.25) {};
\draw[-latex] (1.25,1.25) -- node[below] {$\epsilon$} (0.25,1.25);
\end{tikzpicture}}
\vspace{-18mm}

\begin{itemize}
\setlength\itemsep{4mm}
\item<1-> Consider a metric space $(X,d)$

\item<1-> ``Set of points within distance $\epsilon$ from a point $x$'' \vspace{1mm}
\begin{itemize} 
  \setlength\itemsep{1.5mm}
  \item The \textcolor{blue}{open ball} of radius $\epsilon$ centered at $x$ is \vspace{-1mm} \[\color{blue} B_d(x,\epsilon) \triangleq \left\{ y \in X | d(x,y) < \epsilon \right\} \vspace{-5mm} \]
  \item $P=$``For all $a\!\in\! B_d (x,\epsilon)$, there is $\delta\!>\!0$ s.t. $B_d (a,\delta) \subset B_d (x,\epsilon)$''
\end{itemize}

\item<2-> ``Infinite list $x_1,x_2,x_3,\ldots$ of points in $X$'' \vspace{1mm}
\begin{itemize} 
  \setlength\itemsep{1.5mm}
  \item A \textcolor{blue}{sequence} $x_i \in X$ for $i\in \mathbb{N}$ equivalent to $x_i = f(i)$ for $f:\mathbb{N}\to X$
  \item Ex. For $X=\mathbb{R}$ and $d(x,y)=|x-y|$, let \textcolor{blue}{$x_n = \left(1+\frac{1}{n}\right)^n$} for $n\in \mathbb{N}$
\end{itemize}

\item<3-> ``A sequence of points approaches another point'' \vspace{1mm}
\begin{itemize} 
  \setlength\itemsep{1.5mm}
  \item A sequence $x_n$ \textcolor{blue}{converges} to $x\in X$ (denoted \textcolor{blue}{$x_n \to x$})  if, for any $\epsilon >0$, there is natural number $N$ such that $d(x,x_n) < \epsilon$ for all $n>N$
\end{itemize}
  
\end{itemize}
\end{frame}


\begin{frame}{Convergence: Examples and Counterexamples}

\begin{definition}
A sequence $x_1,x_2,\ldots$ in $(X,d)$ is a \textcolor{blue}{Cauchy sequence} if, for any $\epsilon >0$, there is a natural number $N$ (depending on $\epsilon$) such that, for all $m,n > N$,\vspace{-2mm}
\begin{equation*}
d \left( x_m, x_n \right) < \epsilon
\end{equation*}
\end{definition}

\begin{itemize}
\setlength\itemsep{3mm}
\item<1-> Theorem: Every convergent sequence is a Cauchy sequence \vspace{1mm}
\begin{itemize} 
  \setlength\itemsep{1.5mm}
  \item Prove on board \pause
\end{itemize}
\item Converse? \pause No, there is a counterexample \vspace{1mm}
\begin{itemize} 
  \setlength\itemsep{1.5mm}
  \item Metric space $(X,d)$ with $X=\mathbb{Q}$ and $d(x,y)=|x-y|$
  \item Sequence $x_1 = 2$ and $x_{n+1} = f(x_n) \triangleq \frac{1}{2} x_n + 1/x_n \in \mathbb{Q} $
  \item One can show $x_n$ is a Cauchy sequence and $|x_n - \sqrt{2}|\to 0$
\end{itemize}
\item<4-> But, according to definition \textcolor{blue}{sequence does not converge!}  \vspace{1mm}
\begin{itemize} 
  \setlength\itemsep{1.5mm}
  \item Convergence requires limit lives in $X$ but $\sqrt{2} \notin \mathbb{Q}$
\end{itemize}
\end{itemize}
\end{frame}

  
\begin{frame}{Metric Topology}

A topology is a collection of ``open'' sets satisfying certain properties

\begin{definition}
Let $W$ be a subset of a metric space $(X,d)$.
The set $W$ is called \textcolor{blue}{open} if, for every $w\in W$, there is an $\epsilon>0$ such that $B_d (w,\epsilon) \subseteq W$.
\end{definition}

\begin{definition}
Subset $W$ of $(X,d)$ is \textcolor{blue}{closed} if its complement $W^c = X-W$ is open.
\end{definition}

\begin{theorem}
%For any metric space $(X,d)$,
\begin{enumerate}
\item $\emptyset$ and $X$ are open
\item any union of open sets is open
\item any finite intersection of open sets is open
\end{enumerate}
\end{theorem}

\begin{center}
\scalebox{0.7}{%}
\begin{tikzpicture}[scale=0.8]
\draw[very thick] (-0.5,-1.5) rectangle (3.5,2) node[below left=0.2cm and 0.2cm] {$X$};
\draw [thick,dashed,fill=blue!10] plot [smooth cycle] coordinates {(0,0) (1,1) (3,1) (2,0) (2,-1)} node[above left=0.65cm and 0.15cm] {open};
\end{tikzpicture}}
\hspace{5mm}
\scalebox{0.7}{%}
\begin{tikzpicture}[scale=0.8]
\draw[very thick] (-0.5,-1.5) rectangle (3.5,2) node[below left=0.2cm and 0.2cm] {$X$};
\draw [thick,fill=blue!10] plot [smooth cycle] coordinates {(0.5,0) (1,1.5) (2.5,1) (2.5,0) (1.5,-1)} node[above left=0.75cm and -0.6cm] {closed};
\end{tikzpicture}}
\hspace{5mm}
\scalebox{0.7}{%}
\begin{tikzpicture}[scale=0.8]
\draw[very thick] (-0.5,-1.5) rectangle (3.5,2) node[below left=0.2cm and 0.2cm] {$X$};
%\clip (1,1) circle(0.5);
\fill[blue!10,even odd rule] plot [smooth cycle] coordinates {(0.5,0) (1,1.5) (2.5,1) (2.5,0) (1.5,-1)} (1.5,0.5) circle (0.5);
\draw[thick,dashed] (1.5,0.5) circle (0.5);
\draw [thick] plot [smooth cycle] coordinates {(0.5,0) (1,1.5) (2.5,1) (2.5,0) (1.5,-1)} node[above left=0.35cm and -0.6cm] {neither};
\end{tikzpicture}}
\end{center}  

\end{frame}

\begin{frame}{Interior, Limit points, and Closure}

For a metric space $(X,d)$ and subset $W \subseteq X$:

\begin{definition}
A point $w\in W$ is in the \textcolor{blue}{interior} of $W$ (denoted $\Interior{W}$) if there is a $\delta >0$ such that, for all $x\in X$ with $d(x,w)<\delta$, it follows that $x\in W$.
\end{definition}

\begin{definition}
A point $w\in W$ is a \textcolor{blue}{limit point} of $W$ if there is a sequence of distinct elements, $w_1,w_2,\ldots\in W$, that converges to $w$.
\end{definition}

\begin{definition}
A point $x\in X$ is in the \textcolor{blue}{closure} of $W$ (denoted $\Closure{W}$) if, for all $\delta >0$, there is a $w\in W$ such that $d(x,w)<\delta$.
\end{definition}

\begin{itemize}
\setlength\itemsep{3mm}
\item<1-> Properties \vspace{1mm}
\begin{itemize} 
  \setlength\itemsep{1.5mm}
  \item The interior $\Interior{W}$ is open (see definition)
  \item $W$ is closed if and only if it contains all of its limit points
  \item Closure $\Closure{W}$ equals union of $W$ and all its limit points (thus is closed)
\end{itemize}
\end{itemize}

\end{frame}


\begin{frame}{Continuity}

Let $f \colon X \rightarrow Y$ be a function between metric spaces $(X,d_X)$ and $(Y,d_Y)$:

\begin{definition}
The function $f$ is \textcolor{blue}{continuous at $x_0 \in X$} if, for any $\epsilon> 0$, there exists a $\delta >0$ such that, for all $x\in X$ satisfying $d_X(x_0,x)< \delta$, \vspace{-2mm}
\[ d_Y \left( f(x_0),f(x) \right) < \epsilon .\]
\end{definition}

\begin{theorem}<2-> 
If $f$ is continuous at $x_0$, then $f(x_n) \to f(x_0)$ for all sequences $x_1,x_2,\ldots \in X$ such that $x_n \to x_0$.
Conversely, if $f(x_n) \to f(x_0)$ for all sequences $x_1,x_2,\ldots  \in X$ such that $x_n \to x_0$, then $f$ is continuous at $x_0$.
\end{theorem}

\begin{itemize}
\setlength\itemsep{1.5mm}
\item<3-> $f$ is called \textcolor{blue}{continuous} if it is continuous at all $x_0 \in X$
\item<3-> $f$ is \textcolor{blue}{uniformly continuous} if $\delta$ can be chosen independently of $x_0$
\end{itemize}

\begin{definition}<4->
A function $f  \colon X \rightarrow Y$ is called \textcolor{blue}{Lipschitz continuous} on $A \subseteq X$ if there is a constant $L \in \RealNumbers$ such that $d_Y (f(x),f(y)) \leq L d_X (x,y)$ for all $x,y\in A$.
\end{definition}



\end{frame}


\begin{frame}{Completeness}

\begin{definition}<1->
A metric space $(X,d)$ is said to be \textcolor{blue}{complete} if every Cauchy sequence in $(X,d)$ converges to a limit $x \in X$.
\end{definition}

\begin{example}<1->
Consider the sequence $x_n \in \mathbb{Q}$ defined by $x_1 = 2$ and $\smash{x_{n+1} = \frac{1}{2}x_n + 1/x_n}$.
We have seen that this sequence satisfies $|x_n - \sqrt{2}|\to 0$ but $\sqrt{2}$ is not rational.
Thus, the standard metric space of rationals is not complete.
\end{example}

\begin{definition}<2->
A subset $A$ of a metric space $(X,d)$ is \textcolor{blue}{dense} in $X$ if every $x\in X$ is a limit point of the set $A$.
This is equivalent to the closure $\overline{A}$ being equal to $X$.
\end{definition}

\begin{alertblock}{Key Point}<3->
The standard metric space of real numbers is a complete metric space.
This can be shown using Cauchy sequences of rational numbers because $\mathbb{Q}$ is dense in $\mathbb{R}$.  Note: proof not discussed but available on website.
\end{alertblock}

\end{frame}


\begin{frame}{Contraction Mapping Theorem}

\begin{definition}<1->
Let $A$ be a subset of a metric space $(X,d)$ and $f \colon X \rightarrow X$ be a function.
Then, $f$ is a \textcolor{blue}{contraction} on $A$ if $f(A) \subseteq A$ and there exists a constant $\gamma < 1$ such that $d \left( f(x),f(y) \right) \leq \gamma  d(x,y)$ for all $x,y\in A$.
\end{definition}

\begin{example}
Consider metric space $X=[0,1]$ with absolute distance.
Define $f\colon X \to X$ by $f(x) = 1-\frac{1}{2}x$ and observe $|f(x)-f(y)| = \frac{1}{2}|x-y|$.
\end{example}

\begin{theorem}[Contraction Mapping Theorem]<2->
Let $(X,d)$ be a complete metric space and $f$ be contraction on a closed subset $A \subseteq X$.
Then, $f$ has a unique fixed point $x^*$ in $A$ such that $f(x^*) = x^*$ and $x_{n+1} = f(x_n)$ converges to $x^*$ from any initial $x_1 \in A$.  Moreover, $x_n$ satisfies the error bounds: \\[1mm] \hspace{3mm} $d(x^*,x_n) \leq \gamma^{n-1} d(x^*,x_1)$ and $d(x^*,x_{n+1}) \leq d(x_n,x_{n+1})\gamma /(1-\gamma)$.
\end{theorem}

\end{frame}

\begin{frame}{Applications of the Contraction Mapping Theorem}

The following important results in applied mathematics have relatively simple proofs based on the contraction mapping theorem.
\vspace{2mm}

\begin{itemize}
\setlength\itemsep{3mm}
\item<1-> Picard's uniqueness theorem for differential equations \vspace{1mm}
\begin{itemize} 
  \setlength\itemsep{1.5mm}
  \item Differential equation $y'(t) = f(t,y(t))$ for $t\in [a,b]$ with $y(a)=y_0$ 
  \item  Assume $f(t,y)$ is Lipschitz continuous in $y$ for $t\in [a,b]$
  \item Then, solution $y(t)$ exists and is unique for $t\in [a,b]$
\end{itemize}
\item<2-> Implicit function theorem \vspace{1mm}
\begin{itemize} 
  \setlength\itemsep{1.5mm}
  \item Let $f\colon \mathbb{R}^n \times \mathbb{R}^{m} \to \mathbb{R}^m$ be continuously differentiable on open $A$
  \item Let $g\colon \mathbb{R}^n \to \mathbb{R}^{m}$ be defined implicitly by $f(x,g(x))=0$
  \item For $x_0\in A$, assume $f(x_0,y_0)=0$ and $y$-Jacobian invertible at $(x_0,y_0)$
  \item Then, $g(x)$ exists and is unique in some neighborhood of $x_0$
\end{itemize}
\item<3-> Dynamic Programming for a Markov Decision Process (MDP) \vspace{1mm}
\begin{itemize} 
  \setlength\itemsep{1.5mm}
  \item State-action $(s,a)$ defines probability $p(s'|s,a)$ and reward $R(s,a)$
  \item Finite state + discounted reward $\Rightarrow$ stationary optimal policy
\end{itemize}
\end{itemize}

\end{frame}

\begin{frame}{Contraction Mapping Example}

\begin{minipage}{0.45\textwidth}
Starting from $x_1 = 0.2$, define $x_{n+1}= \cos(x_n)$ and plot the points $(x_n,x_{n+1})$.  Each point is connected to the slope-1 line to emphasize the path taken.
\end{minipage}
\begin{minipage}{0.5\textwidth}
\scalebox{0.75}{% This file was created by matlab2tikz v0.3.0.
% Copyright (c) 2008--2012, Nico Schlömer <nico.schloemer@gmail.com>
% All rights reserved.
% 
% The latest updates can be retrieved from
%   http://www.mathworks.com/matlabcentral/fileexchange/22022-matlab2tikz
% where you can also make suggestions and rate matlab2tikz.
% 
% 
% 
% defining custom colors
\definecolor{mycolor1}{rgb}{0,0.447,0.741}
\definecolor{mycolor2}{rgb}{0.85,0.325,0.098}
\definecolor{mycolor3}{rgb}{1,0,1}
\begin{tikzpicture}
\begin{axis}[%
width=3in,
height=2.5in,
every outer x axis line/.append style={darkgray!60!black},
every x tick label/.append style={font=\color{darkgray!60!black}},
grid=major,
xmin=0, xmax=1,
xlabel={$x_n$},
every outer y axis line/.append style={darkgray!60!black},
every y tick label/.append style={font=\color{darkgray!60!black}},
ymin=0, ymax=1,
ylabel={$x_{n+1}$}]
\addplot [
thick,
color=mycolor1,
solid,
forget plot
]
coordinates{
 (0,0)(0.01,0.01)(0.02,0.02)(0.03,0.03)(0.04,0.04)(0.05,0.05)(0.06,0.06)(0.07,0.07)(0.08,0.08)(0.09,0.09)(0.1,0.1)(0.11,0.11)(0.12,0.12)(0.13,0.13)(0.14,0.14)(0.15,0.15)(0.16,0.16)(0.17,0.17)(0.18,0.18)(0.19,0.19)(0.2,0.2)(0.21,0.21)(0.22,0.22)(0.23,0.23)(0.24,0.24)(0.25,0.25)(0.26,0.26)(0.27,0.27)(0.28,0.28)(0.29,0.29)(0.3,0.3)(0.31,0.31)(0.32,0.32)(0.33,0.33)(0.34,0.34)(0.35,0.35)(0.36,0.36)(0.37,0.37)(0.38,0.38)(0.39,0.39)(0.4,0.4)(0.41,0.41)(0.42,0.42)(0.43,0.43)(0.44,0.44)(0.45,0.45)(0.46,0.46)(0.47,0.47)(0.48,0.48)(0.49,0.49)(0.5,0.5)(0.51,0.51)(0.52,0.52)(0.53,0.53)(0.54,0.54)(0.55,0.55)(0.56,0.56)(0.57,0.57)(0.58,0.58)(0.59,0.59)(0.6,0.6)(0.61,0.61)(0.62,0.62)(0.63,0.63)(0.64,0.64)(0.65,0.65)(0.66,0.66)(0.67,0.67)(0.68,0.68)(0.69,0.69)(0.7,0.7)(0.71,0.71)(0.72,0.72)(0.73,0.73)(0.74,0.74)(0.75,0.75)(0.76,0.76)(0.77,0.77)(0.78,0.78)(0.79,0.79)(0.8,0.8)(0.81,0.81)(0.82,0.82)(0.83,0.83)(0.84,0.84)(0.85,0.85)(0.86,0.86)(0.87,0.87)(0.88,0.88)(0.89,0.89)(0.9,0.9)(0.91,0.91)(0.92,0.92)(0.93,0.93)(0.94,0.94)(0.95,0.95)(0.96,0.96)(0.97,0.97)(0.98,0.98)(0.99,0.99)(1,1) 
};
\addplot [
very thick,
color=red, %mycolor2,
solid,
forget plot
]
coordinates{
 (0,1)(0.01,0.999950000416665)(0.02,0.999800006666578)(0.03,0.999550033748988)(0.04,0.999200106660978)(0.05,0.998750260394966)(0.06,0.998200539935204)(0.07,0.99755100025328)(0.08,0.996801706302619)(0.09,0.995952733011994)(0.1,0.995004165278026)(0.11,0.993956097956697)(0.12,0.992808635853866)(0.13,0.991561893714788)(0.14,0.990215996212637)(0.15,0.988771077936042)(0.16,0.987227283375627)(0.17,0.985584766909561)(0.18,0.983843692788121)(0.19,0.98200423511727)(0.2,0.980066577841242)(0.21,0.978030914724148)(0.22,0.975897449330606)(0.23,0.973666395005375)(0.24,0.97133797485203)(0.25,0.968912421710645)(0.26,0.966389978134513)(0.27,0.963770896365891)(0.28,0.961055438310771)(0.29,0.958243875512697)(0.3,0.955336489125606)(0.31,0.952333569885713)(0.32,0.949235418082441)(0.33,0.946042343528387)(0.34,0.942754665528346)(0.35,0.939372712847379)(0.36,0.935896823677935)(0.37,0.932327345606034)(0.38,0.92866463557651)(0.39,0.924909059857313)(0.4,0.921060994002885)(0.41,0.917120822816605)(0.42,0.913088940312308)(0.43,0.908965749674885)(0.44,0.904751663219963)(0.45,0.900447102352677)(0.46,0.896052497525525)(0.47,0.891568288195329)(0.48,0.886994922779284)(0.49,0.882332858610121)(0.5,0.877582561890373)(0.51,0.872744507645751)(0.52,0.86781917967765)(0.53,0.862807070514761)(0.54,0.857708681363824)(0.55,0.852524522059506)(0.56,0.847255111013416)(0.57,0.841900975162269)(0.58,0.836462649915187)(0.59,0.830940679100164)(0.6,0.825335614909678)(0.61,0.819648017845479)(0.62,0.813878456662534)(0.63,0.808027508312152)(0.64,0.802095757884293)(0.65,0.796083798549056)(0.66,0.789992231497365)(0.67,0.783821665880849)(0.68,0.777572718750928)(0.69,0.771246014997107)(0.7,0.764842187284488)(0.71,0.758361875990508)(0.72,0.751805729140895)(0.73,0.74517440234487)(0.74,0.738468558729588)(0.75,0.731688868873821)(0.76,0.724836010740905)(0.77,0.717910669610943)(0.78,0.710913538012277)(0.79,0.703845315652236)(0.8,0.696706709347165)(0.81,0.689498432951747)(0.82,0.682221207287613)(0.83,0.674875760071267)(0.84,0.667462825841308)(0.85,0.659983145884982)(0.86,0.652437468164052)(0.87,0.644826547240001)(0.88,0.63715114419858)(0.89,0.629412026573697)(0.9,0.621609968270664)(0.91,0.613745749488812)(0.92,0.605820156643463)(0.93,0.597833982287298)(0.94,0.589788025031098)(0.95,0.581683089463883)(0.96,0.573519986072457)(0.97,0.565299531160354)(0.98,0.557022546766217)(0.99,0.548689860581588)(1,0.54030230586814) 
};
\addplot [
thick,
color=mycolor3,
dashed,
forget plot
]
coordinates{
 (0.2,0.2)(0.2,0.980066577841242)(0.980066577841242,0.980066577841242) 
};
\addplot [
thick,
color=mycolor3,
dashed,
forget plot
]
coordinates{
 (0.980066577841242,0.980066577841242)(0.980066577841242,0.556967252809642)(0.556967252809642,0.556967252809642) 
};
\addplot [
thick,
color=mycolor3,
dashed,
forget plot
]
coordinates{
 (0.556967252809642,0.556967252809642)(0.556967252809642,0.848862165658271)(0.848862165658271,0.848862165658271) 
};
\addplot [
thick,
color=mycolor3,
dashed,
forget plot
]
coordinates{
 (0.848862165658271,0.848862165658271)(0.848862165658271,0.660837551116615)(0.660837551116615,0.660837551116615) 
};
\addplot [
thick,
color=mycolor3,
dashed,
forget plot
]
coordinates{
 (0.660837551116615,0.660837551116615)(0.660837551116615,0.789478437766868)(0.789478437766868,0.789478437766868) 
};
\addplot [
thick,
color=mycolor3,
dashed,
forget plot
]
coordinates{
 (0.789478437766868,0.789478437766868)(0.789478437766868,0.704215713341993)(0.704215713341993,0.704215713341993) 
};
\addplot [
thick,
color=mycolor3,
dashed,
forget plot
]
coordinates{
 (0.704215713341993,0.704215713341993)(0.704215713341993,0.762119561760661)(0.762119561760661,0.762119561760661) 
};
\addplot [
thick,
color=mycolor3,
dashed,
forget plot
]
coordinates{
 (0.762119561760661,0.762119561760661)(0.762119561760661,0.723374172105571)(0.723374172105571,0.723374172105571) 
};
\addplot [
thick,
color=mycolor3,
dashed,
forget plot
]
coordinates{
 (0.723374172105571,0.723374172105571)(0.723374172105571,0.749576576331493)(0.749576576331493,0.749576576331493) 
};
\end{axis}
\end{tikzpicture}%}
\end{minipage}
\vspace{1mm}

\begin{itemize}
\setlength\itemsep{1.5mm}
\item<1-> Let $X=[0,1]$ and define $f\colon X\to X$ via $f(x)=\cos(x)$

\item<2-> $\cos([0,1]) = [\cos(1),1]$ because $\cos(x)$ decreasing on $[0,\pi]$

\item<3-> Mean value theorem: $f(y)-f(x) = f'(t) (y-x)$ for some $t\in [x,y]$

\item<4-> $f'(t) = -\sin(t)$ and $\sin([0,1]) = [0,\sin(1)]$ with $\sin(1) \approx 0.84$

\item<5-> $| \cos(y) - \cos(x) \,| \leq 0.85 \, |y-x|$ $\;\Rightarrow$ \textcolor{blue}{$f(x)$ is a contraction on $[0,1]$}

\item<6-> $x_{n+1} \!= \cos(x_n)$  converges to \textcolor{blue}{unique fixed point $x^* \!= \cos(x^*) \!\approx\! 0.739$}
\end{itemize}

\end{frame}  



\begin{frame}{Compactness}

\begin{definition}<1->
A metric space $(X,d)$ is \textcolor{blue}{totally bounded} if, for any $\epsilon > 0$, there exists a finite set of $B_d (x,\epsilon)$ balls that cover (i.e., whose union equals) $X$.
\end{definition}

\begin{definition}<2->
A metric space is \textcolor{blue}{compact} if it is complete~and~totally~bounded.
\end{definition}

\begin{itemize}
\setlength\itemsep{3mm}
\item<3-> Examples \vspace{1mm}
\begin{itemize} 
  \setlength\itemsep{1.5mm}
  \item The closed interval $[0,1]\subset \RealNumbers$ is compact
  \item A subset of Euclidean $\RealNumbers^n$ is compact iff it is closed and bounded
  \item But, the standard metric space of real numbers is not compact because it is not totally bounded.
\end{itemize}

\end{itemize}

\begin{theorem}<4>
A closed subset $A$ of a compact space $X$ is itself a compact space.
\end{theorem}

\end{frame}

\begin{frame}{Compactness and Sequences}

\begin{definition}<1->
Let $x_1,x_2,\ldots \in X$ be a sequence and $n_1,n_2,\ldots\in \mathbb{N}$ be
a strictly increasing sequence.
Then, $x_{n_1},x_{n_2},\ldots$ is called \textcolor{blue}{subsequence}.
\end{definition}

\begin{theorem}<2->
A sequence in a compact metric space has a subsequence that converges.
\end{theorem}

\begin{example}<3->
For the compact metric space $X=[-2,2]\subset \mathbb{R}$ with absolute distance, let $x_n = (-1)^n + \frac{1}{n}$.
Then, subsequence $x_2,x_4,x_6,\ldots$ converges to 1.
\end{example}

\begin{itemize}
\item<3-> Sketch proof on whiteboard in pictures
\end{itemize}

\end{frame}

\begin{frame}{Properties of Real Numbers}

\begin{itemize}
\setlength\itemsep{3mm}
\item<1-> Let us consider extreme values for sets of real numbers \vspace{1mm}

\begin{itemize} 
  \setlength\itemsep{1.5mm}
  \item Extended Real Numbers: $\overline{\RealNumbers} \triangleq \RealNumbers \cup \{ \infty,-\infty\}$
  \item Defines \textcolor{blue}{compact metric space} with metric $d_{\overline{\RealNumbers}} (x,y) \triangleq |\frac{x}{1+|x|} - \frac{y}{1+|y|}|$
  \item ``$x_n \to \infty$'' equivalent to ``$\forall M>0, \, \exists N\in \mathbb{N}, \, \forall n>N, \, x_n > M$''
\end{itemize}
\end{itemize}

\begin{definition}<2->
The \textcolor{blue}{supremum} (or least upper bound) of $X\subseteq \RealNumbers$, denoted \textcolor{blue}{$\sup X$}, is the smallest extended real number $M \in \overline{\RealNumbers}$ such that $x\leq M$ for all $x\in X$.
%It is always well-defined and equals $-\infty$ if $X=\emptyset$.
\end{definition}

\begin{lemma}[supremum sequence]<3->
Let $X$ be a metric space and $f \colon X\rightarrow \RealNumbers$ be a function from $X$ to the real numbers.
Let $M = \sup f(A)$ for some non-empty $A \subseteq X$.
Then, there exists a sequence $x_1,x_2,\ldots \in A$ such that $\lim_n f(x_n) = M$.
\end{lemma}

\begin{itemize}
\item<4-> Sketch proof on whiteboard
\end{itemize}
  
\end{frame}


\begin{frame}{More Properties of Real Numbers}

\begin{definition}<1->
The \textcolor{blue}{maximum} of $X\subseteq \RealNumbers$, denoted \textcolor{blue}{$\max X$}, is the largest value achieved by the set.
It equals the supremum if $\sup X \in X$ and is \textcolor{red}{undefined} otherwise.
\end{definition}

\begin{example}<2->
$X = [1,2) \subset \mathbb{R}$ has $\sup X = 2$ and $\max X$ undefined. \\[1mm]
For $f(x)=\frac{1}{2-x}$, $f(X) = [1,\infty)$ and $\sup f(X) = \infty$.
\end{example}

\begin{itemize}
\item<3-> \textcolor{blue}{infimum}: $\inf X = - \left(\sup \, -\!X\right)$, where $-\!X = \{x\in \mathbb{R}\,|-\!x\in X\}$
\item<3-> \textcolor{blue}{minimum}: $\min X = - \left(\max \, -\!X \right)$, if it exists
\item<3-> supremum and infimum always well-defined but may equal $\pm \infty$
\end{itemize}



\begin{theorem}<4->
Any bounded non-decreasing sequence of real numbers converges to its supremum.
\end{theorem}

\end{frame}  


\begin{frame}{Sequences of Functions}

Let $(X,d_X)$ and $(Y,d_Y)$ be metric spaces and
$f_n \colon X \to Y$ for $n\in \mathbb{N}$ be a sequence of functions mapping $X$ to $Y$.

\begin{definition}
The sequence $f_n$ \textcolor{blue}{converges pointwise} to $f\colon X \to Y$ if, for all $x \in X$, \vspace{-1.5mm}
\[ \lim_{n\to \infty} f_n (x) = f(x) \vspace*{-0.5mm} \]

\end{definition}

\begin{definition}
The sequence $f_n$ \textcolor{blue}{converges uniformly} to $f \colon X \to Y$ if \vspace{-1.5mm}
\[ \forall \epsilon > 0, \, \exists N\in \mathbb{N}, \, \forall n>N, \, \forall x\in X, \, d_Y \left( f_n(x),f(x) \right) < \epsilon. \vspace*{-0.5mm} \]
\end{definition}

\begin{theorem}
If each $f_n$ is continuous and $f_n$ converges uniformly to $f \colon X \to Y$, then $f$ is continuous.
\end{theorem}

\end{frame}

\begin{frame}{Two Important Results}

\begin{theorem}
Let $X$ be a metric space and $f \colon X\rightarrow \RealNumbers$ be a continuous function from $X$ to $\mathbb{R}$.
If $A$ is a compact subset of $X$, then there exists $x\in A$ such that $f(x)=\sup f(A)$ (i.e., $f$ achieves a maximum on $A$).
\end{theorem}

\vspace{3mm}

\begin{theorem}
Let $(X,d)$ be a compact metric space and $C(X)$ be the set of continuous functions mapping $X$ to $\mathbb{R}$.
This set of functions, with the metric \vspace{-1.5mm} \[d_\infty (f,g)=\max_{x\in X} |f(x)-g(x)|, \vspace*{-1.5mm}\] defines a complete metric space.
\end{theorem}

Note: For $f,g\in C(X)$, $d_\infty$ metrizes uniform convergence because \vspace{-1.5mm} \[ \text{``}\max_{x\in X} |f(x)-g(x)| < \epsilon \text{''} \Leftrightarrow \text{``} \forall x\in X, |f(x) - g(x)| < \epsilon \text{''}. \]

\end{frame}

%\begin{frame}{Quantifiers}
%
%\vspace{3mm}
%\hspace*{55mm}
%\input{myfigs/mlp2}
%\vspace{-22mm}
%
%\begin{itemize}
%\item<1-> Neural network
%
%\begin{itemize}
%	\item function $f_\theta$ from $\mathcal{X}=\mathbb{R}^n$ to $\mathcal{Y} = \mathbb{R}^d$ \vspace{1mm}
%	%\item \textcolor{snowymint}{function $f_\theta$} from \textcolor{sail}{$\mathcal{X}=\mathbb{R}^n$} to \textcolor{}{$\mathcal{Y} = \mathbb{R}^d$} \vspace{1mm}
%	%\item $z^{\ell+1} = \sigma ( W_{\ell} z^{\ell} )$
%	\item weights represented by $\theta \in \mathbb{R}^p$ \vspace{1mm}
%\end{itemize}
%
%\vspace{1mm}
%
%\item<2-> Training set  
%
%\begin{itemize}
%	\item Set of tuples $(x,y) \in \mathcal{X} \times \mathcal{Y}$ \vspace{1mm}
%	\item For classification into $d$ classes, let $y\in \mathcal{Y}$ be a one-hot vector \vspace{1mm}
%	\item Entire training set denoted $\mathcal{D} \subset \mathcal{X} \times \mathcal{Y}$ \vspace{1mm}
%\end{itemize}
%
%\only<3->{\vspace{-2.5mm}
%\hspace*{62mm}
%%\includegraphics[width=1.75in]{webfigs/vgg56noshort.png}
%\vspace{-36mm}}
%
%\item<3-> Loss function  
%
%\begin{itemize}
%\vspace{1mm}
%	\item Cross entropy: $L (y,\hat{y}) \triangleq - \sum_{i=1}^d y_i \ln \hat{y}_i$ \vspace{1mm}
%	\item Loss for entire training set is
%	\\[-2mm] \[ \hspace*{-32mm} \mathcal{L}_{\mathcal{D}} (\theta) \triangleq \frac{1}{|\mathcal{D}|}\sum_{(x,y) \in \mathcal{D}} L (y,f_\theta (x)) \]
%	\item note: assume $f_\theta (x) \in \mathbb{R}^d$ non-negative and sums to 1 
%\end{itemize}
%
%
%\end{itemize}
%
%\let\thefootnote\relax\footnotetext{\hspace*{-4mm}  {\tiny Loss landscape taken from \url{https://www.cs.umd.edu/~tomg/img/landscapes/noshort.png} \cite{Li-nips18}}}
%
%\end{frame}
%
%\begin{frame}{Neural Network Training}
%
%\vspace{4mm}
%
%\hspace*{62mm}
%%\includegraphics[width=2in]{myfigs/gradient.pdf}
%\vspace{-46mm}
%
%\begin{itemize}
%
%\item<1-> Loss $\mathcal{L}_\mathcal{D} (\theta)$ \textcolor{blue}{landscape} 
%
%\begin{itemize}
%    \item Adjust $\theta \in \mathbb{R}^p$ to minimize loss \vspace{1mm} 
%    \item Gradient vector $\nabla \mathcal{L}_\mathcal{D} (\theta)$ gives \\ $\theta$ direction of maximum increase \vspace{1mm}
%
%\end{itemize}
%
%\vspace{2mm}
%
%%\only<2->{\hspace*{65mm}
%%\includegraphics[width=2.0in]{webfigs/gradient_descent.jpg}
%%\vspace{-34mm}}
%
%\item<2-> Gradient descent (GD) and SGD
%
%\begin{itemize}
%	%\item Full differential: $\frac{d}{dt} \theta_t = - \nabla \mathcal{L} (\theta_t )$  \vspace{1mm}
%	\item full-batch: $\theta_{t + 1} = \theta_t-\eta \, \nabla \mathcal{L}_\mathcal{D} (\theta_t )$  \vspace{1mm}
%	\item \textcolor{blue}{mini-batch}: subset $\mathcal{S}_t \! \subset \! \mathcal{D}$ with $|\mathcal{S}_t|=m$ \\[-2mm] \[ \hspace*{-35mm} \theta_{t + 1} = \theta_t-\eta \, \nabla \mathcal{L}_{\mathcal{S}_t} (\theta_t ) \]  \vspace{1mm}
%\end{itemize}
%
%\vspace{-5.5mm}
%
%\item<3-> Generalization
%
%\begin{itemize}
%	\item The test set $\mathcal{T} \subset \mathcal{X} \times \mathcal{Y}$ contains held-out training data \vspace{1mm}
%	
%	\item Actual goal is to \textcolor{red}{minimize $\mathcal{L}_\mathcal{T}(\theta)$ without knowing $\mathcal{T}$!} \vspace{1mm}
%	
%\end{itemize}
%
%\end{itemize}
%
%%\let\thefootnote\relax\footnotetext{\hspace*{-4mm} {\tiny Figure: \url{https://reconsider.news/2018/05/09/ai-researchers-allege-machine-learning-alchemy/}} }
%
%\end{frame}



\backupbegin

%\begin{frame}
%\frametitle{Backup Slides}
%\begin{itemize}
%\item Slide numbers not included in denominator!
%\end{itemize}
%\end{frame}

%\begin{frame}[allowframebreaks]
%\frametitle{References}
%\bibliographystyle{alpha}
%\footnotesize
%\bibliography{IEEEabrv,WCLabrv,WCLbib,WCLnewbib}
%\end{frame}

\backupend

\end{document}
